\SourceLeaf{062}
$\frac{\mathrm dF(u,X)}{\mathrm du}$ nancisci valorem $U'$, si statuatur $u=U$, certo esse nequit $U'=0$. Manifesto autem $U'$ erit valor quotientis differentialis partialis $\frac{\mathrm dZ}{\mathrm du}$ pro $u=U,x=X$:
quodsi itaque insuper pro iisdem valoribus ipsarum $u,x$ valorem quotientis differentialis partialis $\frac{\mathrm dZ}{\mathrm dx}$ per $X'$ denotemus, perspicuum est per ea quae in art.
praec. demonstrata sunt, functionem $Z$ per substitutionem
\[
u=U+\frac{XX'}{U'}-\frac{X'x}{U'}
\]
identice evanescere, adeoque per factorem
\[
u+\frac{X'}{U'}x-(U+\frac{XX'}{U'})
\]
indefinite esse divisibilem. Quocirca statuendo $u=xx$, patet, $F(xx,x)$ divisibilem esse per
\[
xx+\frac{X'}{U'}x-(U+\frac{XX'}{U'})
\]
adeoque obtinere valorem 0, si pro $x$ accipiatur radix aequationis
\[
xx+\frac{X'}{U'}x-(U+\frac{XX'}{U'})=0
\]
i. e. si statuatur
\[
x=\frac{-X'\pm\BareRoot{(4UU'U'+4XX'U'+X'X')}}{2U'}
\]
quos valores vel reales esse vel sub forma $g+h\BareRoot{-1}$ contentos constat.

Facile iam demonstratur, per eosdem valores ipsius $x$ etiam functionem $Y$
evanescere debere. Manifesto enim $f(xx,x,\lambda',\lambda'',\lambda'''\text{ etc.})$ est productum ex omnibus $(x-a)(x-b)$ exclusis repetitionibus, et proin $=\upsilon^{m-1}$. Hinc sponte sequitur
\[
\begin{aligned}
f(xx,x,l',l'',l'''\text{ etc.})&=y^{m-1}\\
f(xx,x,L',L'',L'''\text{ etc.})&=Y^{m-1}
\end{aligned}
\]
sive $F(xx,x)=Y^{m-1}$, cuius itaque valor determinatus evanescere nequit, nisi simul evanescat valor ipsius $Y$.

\SecondArticle{20}
Adiumento disquisitionum praecedentium reducta est \textit{solutio} aequationis
$Y=0$, i. e. inventio valoris determinati ipsius $x$, vel realis vel sub forma
$g+h\BareRoot{-1}$ contenti, qui illi satisfaciat, ad solutionem aequationis $F(u,X)=0$,
\SourceLeaf{063}
siquidem determinans functionis $Y$ non fuerit $=0$. Observare convenit, si
omnes coëfficientes in $Y$, i. e. numeri $L',L'',L'''$ etc. sint quantitates reales,
etiam omnes coëfficientes in $F(u,X)$ reales fieri, siquidem, quod licet, pro $X$ quantitas realis accepta fuerit. Ordo aequationis secundariae $F(u,X)=0$ exprimitur per numerum $\frac12m(m-1)$: quoties igitur $m$ est numerus par formae $2^\mu k$ denotante $k$ indefinite numerum imparem, ordo aequationis secundariae exprimitur
per numerum formae $2^{\mu-1}k$.

In casu eo, ubi determinans functionis $Y$ fit $=0$, assignari poterit per
art. 10 functio alia $\mathfrak{Y}$ ipsam metiens, cuius determinans non sit $=0$, et cuius
ordo exprimatur per numerum formae $2^\nu k$, ita ut sit vel $\nu<\mu$, vel $\nu=\mu$. Quaelibet solutio aequationis $\mathfrak{Y}=0$ etiam satisfaciet aequationi $Y=0$: solutio aequationis $\mathfrak{Y}=0$ iterum reducetur ad solutionem alius aequationis, cuius ordo
exprimetur per numerum formae $2^{\nu-1}k$.

Ex his itaque colligimus, generaliter solutionem cuiusvis aequationis, cuius
ordo exprimatur per numerum parem formae $2^\mu k$, reduci posse ad solutionem
alius aequationis, cuius ordo exprimatur per numerum formae $2^{\mu'}k$, ita ut sit
$\mu'<\mu$. Quoties hic numerus etiamnum par est, i. e. $\mu'$ non $=0$, eadem methodus denuo applicabitur, atque ita continuabimus, donec ad aequationem perveniamus, cuius ordo exprimatur per numerum imparem; et huius aequationis
coëfficientes omnes erunt reales, siquidem omnes coëfficientes aequationis primitivae reales fuerunt. Talem vero aequationem ordinis imparis certo solubilem esse
constat, et quidem per radicem realem, unde singulae quoque aequationes antecedentes solubiles erunt, sive per radices reales sive per radices formae $g+h\BareRoot{-1}$.

Evictum est itaque, functionem quamlibet $Y$ formae $x^m-L'x^{m-1}+L''x^{m-2}-$
etc., ubi $L',L''$ etc. sunt quantitates determinatae reales, involvere factorem indefinitum $x-A$, ubi $A$ sit quantitas vel realis vel sub forma $g+h\BareRoot{-1}$ contenta. In casu posteriori facile perspicitur, $Y$ nancisci valorem 0 etiam per substitutionem $x=g-h\BareRoot{-1}$, adeoque etiam divisibilem esse per $x-(g-h\BareRoot{-1})$,
et proin etiam per productum $xx-2gx+gg+hh$. Quaelibet itaque functio
$Y$ certo factorem indefinitum realem primi vel secundi ordinis implicat, et quum
idem iterum de quotiente valeat, manifestum est, $Y$ in factores reales primi vel
secundi ordinis resolvi posse. Quod demonstrare erat propositum huius commentationis.
\begin{center}\rule{3cm}{0.3pt}\end{center}
